\subsection{BOREL SECTIONS}

THEOREM 4. Let X be a Polish space and let R be an equivalence relation on X, such that the equivalence classes mod R are closed in X and the saturation (with respect to R) of each closed set in X is a Borel set. Then there is a Borel set in X which meets each equivalence class in exactly one point.

Consider a metric on X compatible with its topology, and with respect to which X is complete. By Lemma 3 of no. 5, there exists a sifting of X, defined by a sieve \(\mathbf{C} = (\mathbf{C}_{n}, p_{n})\) and a sequence of mappings \((\varphi_{n})\) . For each \(c \in C_{n}\) , let \(g_{n}(c)\) be the saturation of the closed set \(\varphi_{n}(c)\) with respect to R; by hypothesis, \(g_{n}(c)\) is a Borel set in X.

Since each set \(\mathbf{C}_n\) is countable, we can linearly order each \(\mathbf{C}_n\) in such a way that the set of elements smaller than a given element is finite. For each \(c \in \mathbf{C}_n\) we define a set \(h_n(c)\) by induction on \(n\) , as follows. In the first place, for \(c \in \mathbf{C}_0\) , \(h_0(c)\) is the intersection of \(\varphi_0(c)\) and the sets \(\mathbf{X} - g_0(c')\) , where \(c' \in \mathbf{C}_0\) and \(c' < c\) . For \(c \in \mathbf{C}_{n+1}\) , \(h_{n+1}(c)\) is the intersection of \(\varphi_{n+1}(c)\) , \(h_n(\mathbf{P}_n(c))\) and the sets \(\mathbf{X} - g_{n+1}(c')\) for \(c' \in \mathbf{C}_{n+1}\) , \(p_n(c') = p_n(c)\) and \(c' < c\) . The \(h_n(c)\) are clearly Borel sets.

We shall prove the following assertion: for each integer \(n \geqslant 0\) and each equivalence class \(\mathbf{H} \mod \mathbb{R}\) , there is a unique element \(c \in \mathbf{C}_n\) such that \(h_n(c)\) meets \(\mathbf{H}\) , and we have

\[
h (c _ {n}) \cap \mathrm{H} = \varphi_ {n} (c) \cap \mathrm{H},
\]

which is therefore a closed set. For \(n = 0\) , consider the smallest element \(c \in \mathbf{C}_0\) such that \(\varphi_0(c)\) meets \(\mathbf{H}\) ; then \(\varphi_0(c) \cap \mathbf{H}\) does not meet any set \(g_0(c')\) for which \(c' \in \mathbf{C}_0\) and \(c' < c\) ; hence it is contained in \(h_0(c) \cap \mathbf{H}\) and consequently is equal to this set; moreover, we have \(\mathbf{H} \subset g_0(c)\) and therefore \(h_0(c') \cap \mathbf{H}\) is empty for \(c' \in \mathbf{C}_0\) and \(c' > c\) ; thus the assertion is proved for \(n = 0\) . We continue by induction on \(n\) : if there exists \(c \in \mathbf{C}_{n+1}\) such that \(h_{n+1}(c)\) meets \(\mathbf{H}\) , then it follows from the relation \(h_{n+1}(c) \subset h_n(p_n(c))\) and the inductive hypothesis that \(p_n(c)\) is the unique element \(d \in \mathbf{C}_n\) such that \(h_n(d)\) meets \(\mathbf{H}\) . Observe that \(h_n(d)\) , which is contained in \(\varphi_n(d)\) , is contained in the union of the sets \(\varphi_n(c)\) for which \(c \in p_n(d)\) , by the definition of a sifting; there is therefore a smallest element \(c \in p_n(d)\) such that \(\varphi_{n+1}(c)\) meets \(\mathbf{H}\) . We have therefore

\[
\varphi_ {n + 1} (c) \cap \mathrm{H} \subset \varphi_ {n} (d) \cap \mathrm{H} = h _ {n} (d) \cap \mathrm{H}
\]

by the inductive hypothesis. Hence

\[
\varphi_ {n + 1} (c) \cap \mathrm{H} \subset \varphi_ {n + 1} (c) \cap h _ {n} (d),
\]

and since by definition \(\varphi_{n+1}(c) \cap \mathbf{H}\) meets none of the sets \(g_{n+1}(c')\) for which \(c' \in p_n(d)\) and \(c' < c\) , it follows from the definition of \(h_{n+1}(c)\) that \(\varphi_{n+1}(c) \cap \mathbf{H} = h_{n+1}(c) \cap \mathbf{H}\) . Moreover, we have \(\mathbf{H} \subset g_{n+1}(c)\) and therefore, if \(c' \in p_n(d)\) is such that \(c' > c\) , \(h_{n+1}(c') \cap \mathbf{H}\) is empty. Hence the assertion is proved for all \(n\) .

For each integer \(n\) , let \(\mathbf{S}_n\) be the union of the sets \(h_n(c)\) , where \(c\) runs through \(\mathbf{C}_n\) . The set \(\mathbf{S}_n\) is a Borel set, and we have \(\mathbf{S}_{n+1} \subset \mathbf{S}_n\) .

Let S be the intersection of the sets \(S_{n}\) , which is also a Borel set in X; we shall show that S meets each equivalence class H mod R in exactly one point. For each n, let \(c_{n}(H)\) be the unique element \(c \in C_{n}\) such that \(h_{n}(c)\) meets H; then \(\mathbf{S}_{n} \cap \mathbf{H} = \varphi_{n}(c_{n}(\mathbf{H})) \cap \mathbf{H}\) , and \(S \cap H\) is the intersection of the sets \(\varphi_{n}(c_{n}(\mathbf{H})) \cap \mathbf{H}\) . Since the sequence \((c_{n}(\mathbf{H}))\) belongs to L(C), the decreasing sequence of closed sets \(\varphi_{n}(c_{n}(\mathbf{H}))\) , whose diameter tends to o, converges to a point \(x \in X\) , since X is complete. The intersection of the closed sets \(\varphi_{n}(c_{n}(\mathbf{H})) \cap \mathbf{H}\) therefore consists of the point x alone, and the proof of Theorem 4 is complete.

Remark. In particular a closed equivalence relation R satisfies the hypotheses of Theorem 4. When X is a compact metrizable space, Theorem 4 therefore applies to any Hausdorff equivalence relation R, since a Hausdorff equivalence relation on a compact space is closed (Chapter I, § 10, no. 4, Proposition 8).

